\section*{Bab \ref{ch:b2:phylogenetic-trees} --- Pohon Filogeni dan Jam Molekul}

\begin{solution}{exo:b2:phylogenetic-trees:1}
\omterm{def:b2:phylogenetic-trees:tree}{Homologi}: sebuah sifat yang diwarisi dari seorang moyang bersama ---
humerus, radius, dan ulna sayap seekor kelelawar dan sirip seekor paus.
\omterm{def:b2:phylogenetic-trees:tree}{Analogi}: sebuah sifat serupa yang dievolusikan secara bebas --- sayap
kelelawar dan sayap burung sebagai permukaan terbang (tulangnya homolog
sebagai anggota gerak depan, sedangkan sayapnya sebagai sayap tidak),
atau mata kamera gurita dan mata kamera vertebrata. \omterm{def:b2:phylogenetic-trees:tree}{Sinapomorfi}: sebuah
\omterm{def:b2:phylogenetic-trees:tree}{homologi} turunan yang dibagi sebuah \omterm{def:b2:phylogenetic-trees:tree}{klad} dan membatasinya --- bulu bagi
burung, telur amnion bagi amniota.
\end{solution}

\begin{solution}{exo:b2:phylogenetic-trees:2}
Burung: monofiletik. Reptil tanpa burung: parafiletik. Ikan tanpa
tetrapoda: parafiletik (ikan paru lebih dekat kepada kita daripada
kepada ikan trout). Vertebrata yang terbang (kelelawar, burung,
pterosaurus): polifiletik. Mamalia: monofiletik. Prokariota:
parafiletik (bakteri dan arkea, dengan eukariotanya ditinggalkan ---
dan arkea lebih dekat kepada eukariota daripada kepada bakteri).
\end{solution}

\begin{solution}{exo:b2:phylogenetic-trees:3}
Saudara burungnya: buaya. Saudara mamalianya: seluruh amniota lainnya
(kura-kura, kadal, buaya, burung bersama-sama). \omterm{def:b2:phylogenetic-trees:tree}{Klad} terkecil yang
memuat kadal dan burung: (kadal,(buaya, burung)), yaitu diapsid tanpa
kura-kura pada pohon ini. Menukar buaya dan burung pada simpulnya tidak
mengubah apa pun: sebuah simpul dapat diputar dengan bebas.
\end{solution}

\begin{solution}{exo:b2:phylogenetic-trees:4}
Tak berakar: $3$; $3\times 5\times 7 = 105$; $3\times 5\times
7\times 9\times 11\times 13\times 15 = \num{2027025}$. Pohon berakar
bagi 4 takson: tiap pohon dari 3 pohon tak berakarnya memiliki 5
cabang, sehingga memberikan 15.
\end{solution}

\begin{solution}{exo:b2:phylogenetic-trees:5}
$((A,B),(C,D))$: tapak 1, 2, 3 masing-masing satu langkah, tapak 4 dua
langkah ($\mathrm{A}$ dan $\mathrm{G}$ pada kedua sisinya), tapak 5
satu: 6 langkah. $((A,C),(B,D))$: $2 + 2 + 2 + 1 + 1 = 8$.
$((A,D),(B,C))$: $2 + 2 + 2 +
2 + 1 = 9$. Pohon yang pertama paling
parsimonis; padanya tapak 4 bersifat homoplastik (perubahan
$\mathrm{A}\to\mathrm{G}$, atau kebalikannya, terjadi dua kali).
\end{solution}

\begin{solution}{exo:b2:phylogenetic-trees:6}
Gabungkan $A$ dan $B$ pada tinggi 3. Lalu $d(AB,C) = 14$, $d(AB,D) = 16$,
$d(C,D) = 10$: gabungkan $C$ dan $D$ pada tinggi 5. Lalu $d(AB,CD) = (14 + 16
+ 14 + 16)/4 = 15$: akarnya pada tinggi $7.5$. Pohonnya $((A,B),(C,D))$.
\end{solution}

\begin{solution}{exo:b2:phylogenetic-trees:7}
$d = -0.75\ln(1 - 4p/3)$: $0.052$, $0.38$, $1.21$. Pada $p = 0.6$
argumen logaritmanya adalah $0.2$ dan kemiringannya $\mathrm{d}d/\mathrm{d}p
= 1/(1 - 4p/3) = 5$: sebuah galat pencontohan sebesar $0.02$ pada $p$
menjadi $0.1$ pada $d$, dan anggapan laju yang sama pada tiap tapaknya,
yang salah pada tiap gen yang nyata, menambahkan sebuah bias sebesar
itu pula. Gennya jenuh.
\end{solution}

\begin{solution}{exo:b2:phylogenetic-trees:8}
$t = d/2k = 0.02/(2\times 10^{-9}) = 10$ juta tahun. Pada 500 juta
tahun, $d = 2\times 10^{-9}\times 5\times 10^{8} = 1.0$ dan
$p = 0.75(1 - e^{-4/3}) = 0.55$: tiga perempat jalan menuju
kejenuhannya, yang padanya koreksinya curam dan tanggalnya tak
terandalkan --- sebuah gen yang lebih lambat diperlukan bagi pembelahan
semacam itu.
\end{solution}

\begin{solution}{exo:b2:phylogenetic-trees:9}
\qty{52}{\%}: \omterm{def:b2:phylogenetic-trees:tree}{kladnya} diperoleh kembali pada nyaris separuh contoh
ulangnya; datanya nyaris membisu tentangnya, dan pilihan lainnya
bersama-sama sama mungkinnya. \qty{99}{\%}: isyaratnya selaras di
seluruh tapaknya; \omterm{def:b2:phylogenetic-trees:tree}{kladnya} tersokong dengan kuat (mengingat modelnya dan
penjajarannya --- galat sistematis seperti \omterm{prop:b2:phylogenetic-trees:pitfalls}{tarikan cabang panjang} tidak
terdeteksi bootstrapnya). Bagi yang pertama, tambahkan runtunan (lebih
banyak gen), tambahkan takson yang mematahkan cabangnya, lalu periksa
dengan metode yang berbeda.
\end{solution}

\begin{solution}{exo:b2:phylogenetic-trees:10}
Dua garis keturunan yang masing-masing telah berubah pada sebagian
besar tapaknya memiliki, pada tiap tapak yang keduanya berubah
padanya, peluang satu berbanding empat untuk mendarat pada basa yang
sama (dengan empat basa berlaju sama). Jika masing-masing telah berubah
pada \qty{40}{\%} tapaknya, keduanya telah berubah pada \qty{16}{\%}
dan bersesuaian secara kebetulan pada \qty{4}{\%} seluruh tapaknya ---
yang dapat melampaui bagian tapak yang membawa \omterm{def:b2:phylogenetic-trees:tree}{sinapomorfi} sejati yang
menghubungkan masing-masingnya dengan kerabat sejatinya yang lambat
berevolusi. Parsimoni mencacah kesesuaian tanpa menanyakan mengapa ia
terjadi lalu menggabungkan keduanya. Menambahkan takson yang bercabang
sepanjang tiap cabang panjangnya membaginya menjadi ruas yang lebih
pendek, sehingga kesesuaian kebetulannya terbagi di antara beberapa
cabang yang lebih pendek dan \omterm{def:b2:phylogenetic-trees:tree}{sinapomorfi} sejati tiap ruasnya menjadi
terlihat; sebuah model kebolehjadian yang mengharapkan banyak
kesesuaian kebetulan pada cabang yang panjang mengoreksinya secara
langsung.
\end{solution}

\begin{solution}{exo:b2:phylogenetic-trees:11}
$d = -0.75\ln(1 - 0.12) = 0.096$; $t = d/2k = 0.096/(4\times 10^{-8})
= 2.4$ juta tahun. Penggantiannya: $0.096\times \num{16000} =
1500$, dengan ketelitian nisbi $1/\sqrt{1500} = \qty{2.6}{\%}$: sebuah
galat statistik sebesar $\pm 0.06$ juta tahun. Tetapi lajunya
dikalibrasi pada pembelahan lain dan boleh jadi salah dengan faktor dua
bagi garis keturunan ini (keragaman laju antargaris keturunan, dan
selisih antara laju silsilah yang diukur sepanjang beberapa generasi
dan laju yang diukur sepanjang jutaan tahun); dan pemencaran gennya
mendahului pemencaran spesiesnya, karena populasi moyangnya polimorf
--- sehingga pembelahan yang sebenarnya lebih baru daripada pembelahan
gennya, dengan selisih waktu yang bergantung pada ukuran populasi
moyangnya. (Tanggal yang diterima, dari banyak gen inti dan fosil,
adalah 6 sampai 7 juta tahun: jam mitokondria di sini terlalu cepat.)
\end{solution}

\begin{solution}{exo:b2:phylogenetic-trees:12}
Penyortiran garis keturunan yang tak lengkap: ketika dua spesiasinya
berdekatan waktunya, sebuah polimorfisme moyang boleh jadi menyortir
dirinya sehingga pohon sebuah gen mengelompokkan spesiesnya berbeda
dari pohon spesiesnya; pemindahan mendatar: sebuah gen yang diterima
dari garis keturunan lain membawa sejarah garis keturunan itu;
penggandaan gen: membandingkan sebuah paralog pada satu spesies dengan
ortolognya pada spesies lain memberikan tanggal penggandaannya, bukan
tanggal spesiasinya. Sebuah pohon spesies diperoleh dengan memakai
banyak gen dari seluruh genomnya lalu mengambil pohon yang disokong
mayoritasnya (atau sebuah model yang mengharapkan sebagian gen yang
tidak selaras yang diketahui), dengan memilih ortolognya secara
saksama, dan dengan lebih menyukai, pada prokariota, gen ribosom dan
gen informasional yang jarang dipindahkan.
\end{solution}

\begin{solution}{pb:b2:phylogenetic-trees:1}
\textbf{1.} Tapak 1, 2, dan 4 (keadaan yang dibagi $W$ dan $H$ saja,
yang turunan terhadap untanya): \omterm{def:b2:phylogenetic-trees:tree}{sinapomorfi} $\{W,H\}$. Tapak 3
dan 6: \omterm{def:b2:phylogenetic-trees:tree}{sinapomorfi} $\{W,H,C\}$. Tapak 5 adalah sebuah autapomorfi
$W$, yaitu sebuah perubahan pada satu garis keturunan yang tidak
mengelompokkan apa pun; tidak ada tapak yang bersifat homoplastik pada
pohon molekulnya.
\textbf{2.} Satu perubahan tiap tapaknya: 6 langkah.
\textbf{3.} Tapak 1, 2, dan 4 kini masing-masing memerlukan dua
perubahan, tapak 3, 5, dan 6 masing-masing satu: 9 langkah.
\textbf{4.} Pohon paus--kuda nilnya, dengan selisih tiga langkah. Tiga
tapak dari seribu adalah selisih yang kecil yang dapat dijungkalkan
beberapa homoplasi; menguatkannya berarti lebih banyak gen, lebih
banyak takson (ruminansia lainnya, pekari, rusa), dan sebuah bootstrap.
\textbf{5.} \omterm{def:b2:phylogenetic-trees:tree}{Kladnya} muncul lagi pada lima contoh ulang dari enam:
sokongan sedang, jauh di atas kebetulan tetapi kurang dari \qty{95}{\%}
yang secara lazim disebut kuat.
\textbf{6.} \omterm{meth:b2:phylogenetic-trees:parsimony}{Kelompok luarnya} harus terletak di luar kelompoknya tetapi
cukup dekat sehingga runtunannya masih dapat dijajarkan dan belum
jenuh; untanya adalah hewan berkuku genap di luar kelompok
sapi--babi--kuda nil--paus. Seekor ikan akan berbeda pada aras yang
nyaris jenuh, keadaannya tidak akan informatif tentang keadaan mana
yang merupakan keadaan moyang di dalam mamalianya, dan cabang panjangnya
akan menarik garis keturunan dalam yang paling cepat berevolusi.
\textbf{7.} $d = 0.0625$, $0.107$, $0.131$, $0.155$.
\textbf{8.} Gabungkan $W$ dan $H$ pada tinggi $0.031$. Lalu $d(WH,C) =
0.107$, $d(WH,P) = 0.131$, $d(C,P) = 0.131$, dan semuanya ke $M$
$0.155$: gabungkan $WH$ dan $C$ pada $0.054$. Lalu $d(WHC,P) = 0.131$:
gabungkan pada $0.065$. Akhirnya gabungkan $M$ pada $0.078$. Pohonnya
$(M,(P,(C,(W,H))))$.
\textbf{9.} Ya: pohon yang sama, persarangan yang sama.
\textbf{10.} Seekor paus yang telah berubah dua kali lebih cepat akan
lebih jauh dari segalanya, kuda nil termasuk; UPGMA, yang menaruh tiap
ujungnya pada tinggi nol lalu menggabungkan pasangan yang terdekat,
akan menggabungkan kuda nil dengan sapi lebih dahulu lalu meninggalkan
pausnya di luar --- tepat pohon para ahli anatomi, karena alasan yang
salah.
\textbf{11.} $0.0625\times 1000 = 62.5$ penggantian; simpangan baku
Poissonnya $\sqrt{62.5} = 7.9$.
\textbf{12.} $7.9/62.5 = \qty{13}{\%}$.
\textbf{13.} $k = d/2t = 0.107/(1.2\times 10^{8}) = 8.9\times 10^{-10}$
tiap tapak tiap tahun.
\textbf{14.} $t = 0.0625/(2\times 8.9\times 10^{-10}) = 35$ juta
tahun.
\textbf{15.} Babi: $0.131/(1.79\times 10^{-9}) = 73$ juta tahun;
unta: $0.155/(1.79\times 10^{-9}) = 87$ juta tahun.
\textbf{16.} $35 \pm 4.5$ juta tahun (\qty{13}{\%}).
\textbf{17.} Lajunya berskala sebagai $1/t_{\text{kal}}$ dan tanggalnya
sebagai $t_{\text{kal}}$: $\pm 5/60 = \qty{8}{\%}$, sehingga $\pm 3$
juta tahun --- jika digabung dengan galat Poissonnya, kira-kira $\pm 5$
juta tahun.
\textbf{18.} $d = -0.75\ln(1 - 0.6) = 0.69$: pada dasarnya gennya telah
berubah pada dua pertiga tapaknya, koreksinya telah melipatkan $p$
dengan $1.5$ dan kemiringannya $2.5$; ia dekat kejenuhan bagi
pembelahan ini dan tanggalnya bernilai kecil.
\textbf{19.} Ketiadaan sebuah sifat bukanlah sebuah keadaan turunan
yang dibagi: paus kehilangan anggota gerak belakangnya beserta
pergelangan kakinya sama sekali, sehingga tidak ada yang dapat
dikatakan tentang pergelangan kaki mana yang akan ia miliki. Sebuah
\omterm{def:b2:phylogenetic-trees:tree}{sinapomorfi} adalah bukti sebuah \omterm{def:b2:phylogenetic-trees:tree}{klad}; hilangnya pada satu garis
keturunan bukanlah bukti yang melawan keanggotaannya.
\textbf{20.} Fosil pausnya memiliki pergelangan kaki berkatrol ganda:
paus berasal dari seorang moyang yang memilikinya, dan hal itu
menaruhnya di dalam hewan berkuku genap. Ramalan molekulnya dibenarkan
batuannya.
\textbf{21.} Parafiletik: seorang moyang beserta seluruh keturunannya
kecuali garis pausnya.
\textbf{22.} Penyortiran garis keturunan yang tak lengkap atas sebuah
polimorfisme moyang melintasi dua spesiasi yang cepat (sapi, kuda nil,
dan paus terbelah dalam beberapa juta tahun), atau perbandingan
paralog. Putuskan dengan meruntunkan banyak gen lalu mengambil pohon
yang disokong sebagian besarnya, dan gabungannya; periksalah
penggandaan pada tiap keluarganya.
\textbf{23.} $d = -0.75\ln(1 - 0.333) = 0.30$; $k = d/2t =
0.30/(7\times 10^{7}) = 4.3\times 10^{-9}$ tiap tapak tiap tahun, yaitu
lima kali gen intinya.
\textbf{24.} DNA mitokondria digandakan sebuah polimerase yang
penyuntingannya lebih lemah, terpapar radikal oksigen rantai
pernapasannya, dan tidak diperbaiki mesin intinya; ia juga tidak
memiliki rekombinasi, sehingga kerusakannya menimbun.
\textbf{25.} Pohonnya $(M,(P,(C,(W,H))))$: paus adalah kelompok saudara
kuda nil; $k = 8.9\times 10^{-10}$ tiap tapak tiap tahun; pemencaran
paus--kuda nilnya $35 \pm 5$ juta tahun.
\end{solution}
